\documentclass[main.tex]{subfiles}


\begin{document}

%from pagenumber to up
% \phantomsection 
% \hypertarget{toc}{}

 

 

 
   

\section{Работа газа}


Пусть сосуд с газом, находящимся под неподвижным гладким массивным поршнем, поставили на огонь (рис.\,\ref{fig:p82}, слева).

    \begin{figure}[H]
    \centering

    \begin{tikzpicture}[font=\small,            line cap=round,                             >={Stealth[inset=0pt,round,length=3mm, angle'=20]},vec/.style={>={Stealth[inset=0pt,round,length=3.5mm, angle'=20]}}]


% plane
\filldraw[lightgray,semitransparent] (0,0) -- (5,0) --++ (0,-.3) --++ (-5,0) -- cycle;
\draw[] (0,0) -- (5,0);


% %solid
% \fill[lightgray] (2.5,1) rectangle++ (2,.2);
%     % red part
%     \fill[red,path fading=east] (2.5,1) rectangle++ (1,.2);
%     \draw[] (2.5,1) rectangle++ (2,.2);


%plunger
\fill[lightgray] (1.5,2.6) --++ (2,0) --++ (0,.5) --++ (-2,0) -- cycle;
\draw[] (1.5,2.6) --++ (2,0) ++ (0,.5) --++ (-2,0);
% \node[above] at (2.5,3.1) {$M$};


%cyl
\draw (1.5,1.1) coordinate (gl) ++ (0,3.5) --++ (0,-3.5) --++ (2,0) coordinate (gr) --++ (0,3.5);



%sup
\fill[lightgray,semitransparent] (gl) ++ (.25,0) --++ (-.5,0) --++ (0,-1.1) -- cycle;
\draw[] (gl) ++ (.25,0) --++ (-.5,0) --++ (0,-1.1) -- cycle;

\fill[lightgray,semitransparent] (gr) ++ (-.25,0) --++ (.5,0) --++ (0,-1.1) -- cycle;
\draw[] (gr) ++ (-.25,0) --++ (.5,0) --++ (0,-1.1) -- cycle;

%gas
\foreach \i in {1,...,100}
  \fill[Green,shift={(1.54,1.14)}] ({rnd*1.92cm}, {rnd*1.42cm}) circle (.02);
\node[fill=white] at (2.5,1.85) {$P,V_1$};

%flame

\begin{scope}[yshift=.1cm]
    
\fill[red,nearly opaque] (2.5,0) 
to [out=180,in=-90] (2.2,.3) 
to [out=90,in=-90] (2.5,1) 
to [out=-45,in=90] (2.7,.5) 
to [out=45,in=-90] (2.75,.7) 
to [out=-45,in=0] (2.5,0);

\fill[scale around={.8:(2.5,0)},orange] (2.5,0) 
to [out=180,in=-90] (2.2,.3) 
to [out=90,in=-90] (2.55,1) 
to [out=-45,in=90] (2.65,.4) 
to [out=45,in=-90] (2.8,.55) 
to [out=-45,in=0] (2.5,0);

\fill[scale around={.5:(2.5,0)},yellow] (2.5,0) 
to [out=180,in=-90] (2.2,.3) 
to [out=90,in=-110] (2.55,1) 
to [out=-70,in=150] (2.8,.5) 
to [out=-100,in=90] (2.8,0.3)
to [out=-90,in=0] (2.5,0);

\end{scope}

\filldraw[] (2,0) rectangle++ (1,.1);

%%%%%%%%%%%%%%%%%%%%%%%%%%%%%%%%%%%%%%%
%right pic

\begin{scope}[xshift=7cm]

% plane
\filldraw[lightgray,semitransparent] (0,0) -- (5,0) --++ (0,-.3) --++ (-5,0) -- cycle;
\draw[] (0,0) -- (5,0);


% %solid
% \fill[lightgray] (2.5,1) rectangle++ (2,.2);
%     % red part
%     \fill[red,path fading=east] (2.5,1) rectangle++ (1,.2);
%     \draw[] (2.5,1) rectangle++ (2,.2);


%plunger
\begin{scope}[yshift=1cm]
    
\fill[lightgray] (1.5,2.6) --++ (2,0) --++ (0,.5) --++ (-2,0) -- cycle;
\draw[] (1.5,2.6) --++ (2,0) ++ (0,.5) --++ (-2,0);
% \node[above] at (2.5,3.1) {$M$};
\end{scope}


%cyl
\draw (1.5,1.1) coordinate (gl) ++ (0,3.5) --++ (0,-3.5) --++ (2,0) coordinate (gr) --++ (0,3.5);



%sup
\fill[lightgray,semitransparent] (gl) ++ (.25,0) --++ (-.5,0) --++ (0,-1.1) -- cycle;
\draw[] (gl) ++ (.25,0) --++ (-.5,0) --++ (0,-1.1) -- cycle;

\fill[lightgray,semitransparent] (gr) ++ (-.25,0) --++ (.5,0) --++ (0,-1.1) -- cycle;
\draw[] (gr) ++ (-.25,0) --++ (.5,0) --++ (0,-1.1) -- cycle;

%gas
\foreach \i in {1,...,100}
  \fill[Green,shift={(1.54,1.14)}] ({rnd*1.92cm}, {rnd*2.42cm}) circle (.02);
\node[fill=white] at (2.5,2.35) {$P,V_2$};

%flame

\begin{scope}[yshift=.1cm]
    
\fill[red,nearly opaque] (2.5,0) 
to [out=180,in=-90] (2.2,.3) 
to [out=90,in=-90] (2.5,1) 
to [out=-45,in=90] (2.7,.5) 
to [out=45,in=-90] (2.75,.7) 
to [out=-45,in=0] (2.5,0);

\fill[scale around={.8:(2.5,0)},orange] (2.5,0) 
to [out=180,in=-90] (2.2,.3) 
to [out=90,in=-90] (2.55,1) 
to [out=-45,in=90] (2.65,.4) 
to [out=45,in=-90] (2.8,.55) 
to [out=-45,in=0] (2.5,0);

\fill[scale around={.5:(2.5,0)},yellow] (2.5,0) 
to [out=180,in=-90] (2.2,.3) 
to [out=90,in=-110] (2.55,1) 
to [out=-70,in=150] (2.8,.5) 
to [out=-100,in=90] (2.8,0.3)
to [out=-90,in=0] (2.5,0);

\end{scope}

\filldraw[] (2,0) rectangle++ (1,.1);

\end{scope}


%\Longrightarrow
\path (6,2.3) node {$\Longrightarrow$};


    \end{tikzpicture}
\caption{Расширение газа}
\label{fig:p82}
\end{figure}

Газ начального объема~$V_1$ медленно получает тепло от огня, так что давление~$P$ газа не меняется~--- оно равно давлению почти покоящегося поршня на газ (рис.\,\ref{fig:p82}, слева). При постоянном давлении температура газа медленно повышается, то есть его молекулы начинают двигаться быстрее~--- они сильнее толкают поршень, и он постепенно поднимается (рис.\,\ref{fig:p82}, справа). За время наблюдения газ расширяется до объема~$V_2$, совершая работу.

\textbf{Работа газа} ($A$ [Дж])~--- это работа сил давления газа. При \emph{постоянном давлении} газа ее находят по формуле:
\begin{equation}
    A=P\Delta V,
\end{equation}
где $P$~--- давление, а $\Delta V$~--- изменение объема газа.

При расширении газ совершает положительную работу ($A>0$), при сжатии~--- отрицательную ($A<0$). Так, в рассмотренном примере (рис.\,\ref{fig:p82}) газ совершает положительную работу $A=P(V_2-V_1)$.

Описанный выше процесс c газом можно изобразить на $PV$"=диаграмме (рис.\,\ref{fig:p83}).

\begin{figure}[H]
    \centering

    \begin{tikzpicture}[font=\small,            line cap=round,                             >={Stealth[inset=0pt,round,length=3mm, angle'=20]},vec/.style={>={Stealth[inset=0pt,round,length=3.5mm, angle'=20]}}]


\draw[->] (-.3,0) --++ (0:4.3cm) node[below,black] {$V$};
\draw[->] (0,-.3) --++ (90:3.3cm) node[left,black] {$P$};
\node[below left] at (0,0) {$O$};


%S
\fill[red,nearly transparent] (0,2) ++ (1,0) --++ (2,0) --++ (0,-2) --++ (-2,0) -- cycle;
\node at (2,1) {$A$};

%red
\draw[densely dashed] (0,2) node[left] {$P$} --++ (0:1cm) coordinate (v1);
\draw[thick,red] (v1) --++ (0:2cm);

\draw[densely dashed] (v1) --++ (-90:2cm) node[below] {$V_1$};
\draw[densely dashed] (v1) ++ (0:2cm) --++ (-90:2cm) node[below] {$V_2$};



    \end{tikzpicture}
\caption{Работа газа как площадь}
\label{fig:p83}
\end{figure}

Как видно, \emph{работа газа} равна \emph{площади фигуры} (взятой с соответствующим знаком) между графиком и осью~$V$ на $PV$"=диаграмме для соответствующего изменения объема. 

\textbf{Работа внешних сил}~$A'$ над газом связана с работой газа~$A$ так:
\begin{equation}
    A'=-A.
\end{equation}













 
\end{document}